% LTex: language=cs
\def\Abstrakt{%
\begin{otherlanguage}{czech}
Tato práce se zabývá problémem dispečerského řízení vlaků (TDP) a návrhem heuristiky schopné rozhodování v reálném čase pro komplexní železniční sítě. Naše metoda využívá přístup shora dolů pro rozhodování o trasování, motivovaný strukturní analýzou instancí, včetně bezpečnostních hledisek, a formulací pomocí cest a cyklů pro řešení konfliktů, což odstraňuje nutnost používat omezení s velkým M.

Procedury pro trasování i řešení konfliktů vycházejí z efektivního agregačního postupu a inovativní struktury grafu propojení (link graph). Agregační postup zmenšuje velikost problému a lze jej využít pro jakoukoli rozvrhovací úlohu s alokací zdrojů. Graf propojení uchovává informace o závislostech zdrojů pro aktuální výběr tras a umožňuje nám vytvářet strukturní implikační pravidla, která lze použít k identifikaci kritických částí při rozhodování o trasování a ke snížení počtu požadovaných proměnných a omezení při řešení konfliktů.

Heuristika byla použita k vyřešení většiny instancí v knihovně benchmarků DISPLIB a přinesla výrazné zlepšení rychlosti ve srovnání s vítěznými řešeními, i když za cenu občasných suboptimálních řešení. Výsledky jasně ukazují výhody a nevýhody formulace cest a cyklů a ilustrují, jak se struktura grafu propojení využívá v praxi.
\end{otherlanguage}
}
